\chapter{Produit des nombres relatifs}

\begin{itemize}
\item Effectuer : 
\begin{itemize}
\it $(+17)(-13)$.
\it $\left(+\frac23\right)\left(+\frac57\right)$.
\it $(-0,7)\left(+\frac7{11}\right)$.
\it $(-14)(-11)$.
\it $\left(-\frac{11}{11}\right)\left(-\frac34\right)$.
\it $\left(-\frac32\right)(-0,5)$.
\it $(+12)(+4)$.
\it $\left(+\frac59\right)\left(-\frac75\right)$.
\it $(+0,7)\left(-\frac49\right)$.
\end{itemize}
\item Effectuer 
\begin{itemize}
\it $(-3)(+7)(-5)(-4)$.
\it $(+1,7)(-2,5)(-3)(-1,9)$.
\it $(-1)(+1)(-1)(+1)(-1)(-1)$.
\it $\left(-\frac35\right)(+3)\left(-\frac53\right)$.
\it $\left(+\frac12\right)\left(-\frac14\right)\left(-\frac48\right)$.
\it $\left(+\frac12\right)\left(-\frac6{11}\right)\left(-\frac78\right)\left(-\frac{11}{15}\right)$.
\end{itemize}
\item Calculer de deux façons différentes les produits suivants : 
\begin{itemize}
\it $(-5+7-9-1)(-3)$.
\it $(+12-17-9)(+7)$.
\it $(-12+0,7)(+0,5)$.
\it $(-13+11,7+12,5)(-4,7)$.
\it $\left(-\frac12+\frac23+\frac35\right)(-4)$.
\it $\left(-\frac35+\frac79+\frac{11}{13}\right)\left(+\frac23\right)$.
\end{itemize}
\item Calculer de deux façons différentes les produits suivants : 
\begin{itemize}
\it $(-5+12-7)(5-7)$.
\it $\left(\frac35-\frac4{15}\right)\left(\frac27-\frac3{14}\right)$.
\it $(+5+14-13)(-7-2)$.
\it $\left(\frac59-\frac{7}{27}\right)\left(1-\frac35\right)$.
\it $(-2,5+2,7+4,9)(0,75-1)$.
\it $\left(\frac{11}{12}-1\right)\left(\frac{11}{12}+1\right)$.
\end{itemize}
\item Effectuer le plus rapidement possible les opérations suivantes : 
\begin{itemize}
\it $(5-3+2)(12-9)+(15-17)(7-10)$.
\it $(7-17)(10-12) - (15-17)(-14+11)$.
\it $(2,5+0,7)(7,9-8)-0,5(11-8)$.
\it $\left(\frac23+\frac35-1\right)\left(1-\frac79\right)+\frac45\left(17-\frac13\right)$.
\it $[(5-11)-(17-14)][12-(14-11)] + 13(17-12)$.
\it $[14-(0,7-1)]\left[0,5+\left(\frac12-3\right)\right] - [12-(2,4+4,1)][(1,7+2,3)-1]$. 
\end{itemize}
\it La vitesse d'un point mobile sur un axe est $+4$ cm par seconde, et ce mobile passe au point $O$, origine des abscisses au temps $t=0$. Quelles sont les abscisses de ce point aux temps $+3$ et $+7$ ? Quelle est la valeur algébrique du vecteur joignant les $2$ positions de ce point aux temps considérés ? Même question si la vitesse est $-4$ cm par seconde. 
\it Au temps $+5$, l'abscisse d'un mobile est $-16$ cm, au temps $-3$ elle est $+8$. Quelle est la vitesse de ce mobile sachant qu'il est animé d'un mouvement uniforme ? 
\end{itemize}